%%%%%%%%%%%%%%%%%%%%%%%%%%%%%%%%%%%%%%%%%%%%%%%%%%%%%%%%%%%%%%%%%%%%%%
% Lớp 11 | Hàm số mũ và hàm số lôgarít | Phép tính luỹ thừa
% Dạng 1D6?1-4: So sánh các lũy thừa
% Số câu: 24  (sinh tự động bởi tach_cau.py; từ nay thầy sửa trực tiếp file này)
%%%%%%%%%%%%%%%%%%%%%%%%%%%%%%%%%%%%%%%%%%%%%%%%%%%%%%%%%%%%%%%%%%%%%%

%%%%%%%%%%%%%%%% TRẮC NGHIỆM 4 PHƯƠNG ÁN %%%%%%%%%%%%%%%%

% nguồn: Ok_Toan11_TracNghiem_Lan2
\begin{ex}%[1D6N1-4]%[ma:57419afc]
	Khẳng định nào dưới đây  là đúng?
	 \choice
		 {$\left(\dfrac{3}{7}\right)^{\sqrt{3}}>\left(\dfrac{5}{8}\right)^{\sqrt{3}}$}
		 {\True $\left(\dfrac{1}{2}\right)^{-\pi}<\left(\dfrac{1}{3}\right)^{-\pi}$}
		 {$3^{-\sqrt{2}}<\left(\dfrac{1}{5}\right)^{\sqrt{2}}$}
		 {$\left(\dfrac{1}{4}\right)^{-50}<\left(\sqrt{2}\right)^{100}$}
	 \loigiai{
		Khẳng định \lq\lq$\left(\dfrac{1}{2}\right)^{-\pi}<\left(\dfrac{1}{3}\right)^{-\pi}$\rq\rq \, đúng, vì  $\dfrac{1}{2}>\dfrac{1}{3}$.
	} 
\end{ex}

% nguồn: Ok_Toan11_TracNghiem_Lan2
\begin{ex}%[1D6N1-4]%[ma:ba30abe9]
	 Nếu $a^{\tfrac{1}{3}}>a^{\tfrac{1}{6}}$ và $b^{\sqrt{3}}>b^{\sqrt{5}}$ thì
	\choice 
		 {$a<1$; $0<b<1$}
		 {$a>1$; $b<1$}
		{$0<a<1$;  $b<1$}
		{\True $a>1$; $0<b<1$} 
	 \loigiai{
		Ta có $\dfrac{1}{3}>\dfrac{1}{6}$, mà $a^{\tfrac{1}{3}}>a^{\tfrac{1}{6}}$  nên $a>1$.\\
		Lại có  $\sqrt{3}<\sqrt{5}$, mà $b^{\sqrt{3}}>b^{\sqrt{5}}$ nên $0<b<1$.
	} 
\end{ex}

% nguồn: Ok_Toan11_TracNghiem_Lan2
\begin{ex}%[1D6N1-4]%[ma:42676224]
	Tìm tập tất cả các giá trị của $a$  để $\sqrt[21]{a^5}>\sqrt[7]{a^2}$?
	 \choice
		 {$a>0$}
		 {\True $0<a<1$}
		{$a>1$} 
		 {$\dfrac{5}{21}<a<\dfrac{2}{7}$}
	\loigiai{ 
		Ta  có $\sqrt[7]{a^2}=\sqrt[21]{a^6}$.\\
		Do đó $\sqrt[21]{a^5}>\sqrt[7]{a^2} \Leftrightarrow  \sqrt[21]{a^5}>\sqrt[21]{a^6}$, mà $5<6$.\\
		Vậy $0<a<1$. 
	 }
\end{ex}

% nguồn: Ok_TracNghiemLop11_Lan1
\begin{ex}%[1D6N1-4]%[ma:15159a67]
	Cho $\alpha, \beta$ là hai số thực thỏa mãn $2024^{\alpha} > 2024^{\beta}$. Khẳng định nào sau đây  là đúng?
	\choice 
		{\True  $\alpha > \beta$}
		 {$\alpha=\beta$}
		{$\alpha \leq \beta$} 
		{$\alpha <  \beta$}
	\loigiai{ 
		Vì cơ số $a=2024>1$  nên $2024^{\alpha} > 2024^{\beta}\Leftrightarrow \alpha>\beta$.
	 }
\end{ex}

% nguồn: Ok_TracNghiemLop11_Lan1
\begin{ex}%[1D6N1-4]%[ma:03cdc60f]
	Cho số thực dương $a$ thỏa mãn $\left( \dfrac{1}{a}\right)^{24}>\left( \dfrac{1}{a}\right)^{23} $. Khẳng  định nào sau đây đúng?
	\choice 
		{\True  $0<a<1$}
		{$a<\dfrac{1}{2}$} 
		 {$\dfrac{1}{2}<a<1$}
		 {$a>1$}
	 \loigiai{
		Với $a>0$ và $\left( \dfrac{1}{a}\right)^{24}>\left( \dfrac{1}{a}\right)^{23} $, ta suy  ra $\dfrac{1}{a}>1$ hay $a<1$.\\
		Vậy  $0<a<1$.
	 }
\end{ex}

% nguồn: Ok_Toan11_TracNghiem_Lan2
\begin{ex}%[1D6H1-4]%[ma:8d715104]
	Khẳng định nào sau đây đúng? 
	\choice 
		 {$\left(\sqrt{5}+2\right)^{-2\,017}<\left(\sqrt{5}+2\right)^{-2\,018}$}
		{$\left(\sqrt{5}+2\right)^{2\,018}>\left(\sqrt{5}+2\right)^{2\,019}$} 
		{\True $\left(\sqrt{5}-2\right)^{2\,018}>\left(\sqrt{5}-2\right)^{2\,019}$} 
		{$\left(\sqrt{5}-2\right)^{2\,018}<\left(\sqrt{5}-2\right)^{2\,019}$} 
	\loigiai{ 
		Khẳng định  \lq\lq$\left(\sqrt{5}-2\right)^{2\,018}>\left(\sqrt{5}-2\right)^{2\,019}$\rq\rq \, đúng, vì $\heva{& \sqrt{5}-2\in (0;1)\\& 2\,018<2\,019.}$
	 }
\end{ex}

% nguồn: Ok_Toan11_TracNghiem_Lan2
\begin{ex}%[1D6H1-4]%[ma:fcfb8ca6]
	Cho $a=3^{\sqrt{5}}$, $b=3^2$ và $c=3^{\sqrt{6}}$. Mệnh đề nào dưới đây  đúng?
	\choice 
		 {$a<c<b$}
		 {$a<b<c$}
		{\True $b<a<c$} 
		{$c<a<b$} 
	\loigiai{ 
		Vì $2<\sqrt{5}<\sqrt{6}$ nên  $3^2<3^{\sqrt{5}}<3^{\sqrt{6}}$ hay $b<a<c$.
	} 
\end{ex}

% nguồn: Ok_Toan11_TracNghiem_Lan2
\begin{ex}%[1D6H1-4]%[ma:462794ad]
	Cho $a>1$. Mệnh đề nào dưới đây đúng? 
	\choice 
		{$\dfrac{1}{a^{2\,016}}<\dfrac{1}{a^2\,017}$} 
		{$a^{\tfrac{1}{3}}>\sqrt{a}$} 
		{\True $a^{-\sqrt{3}}>\dfrac{1}{a^{\sqrt{5}}}$} 
		{$\dfrac{\sqrt[3]{a^2}}{a}>1$} 
	 \loigiai{
		Vì $a>1$ nên  $a^{\sqrt{3}}<a^{\sqrt{5}}$.\\
		Suy  ra $\dfrac{1}{a^{\sqrt{3}}}>\dfrac{1}{a^{\sqrt{5}}} \Leftrightarrow a^{-\sqrt{3}}>\dfrac{1}{a^{\sqrt{5}}}$.
	} 
\end{ex}

% nguồn: Ok_Toan11_TracNghiem_Lan2
\begin{ex}%[1D6H1-4]%[ma:9414f8ed]
	Trong các mệnh đề sau, mệnh  đề nào \textbf{sai}?
	\choice 
		{\True  $\left(\sqrt{3}-1\right)^{2\,018}>\left(\sqrt{3}-1\right)^{2\,017}$}
		 {$2^{\sqrt{2}+1}>2^{\sqrt{3}}$}
		 {$\left(\sqrt{3}-1\right)^{2\,017}>\left(\sqrt{3}-1\right)^{2\,018}$}
		{$\left(1-\dfrac{\sqrt{2}}{2}\right)^{2\,019}<\left(1-\dfrac{\sqrt{2}}{2}\right)^{2\,018}$} 
	 \loigiai{
		Mệnh đề \lq\lq$\left(\sqrt{3}-1\right)^{2\,018}>\left(\sqrt{3}-1\right)^{2\,017}$\rq\rq \,  sai, vì
		$$\heva{& \sqrt{3}-1\in  (0;1)\\& 2\,018>2\,017} \Rightarrow \left(\sqrt{3}-1\right)^{2\,018}<\left(\sqrt{3}-1\right)^{2\,017}.$$
	} 
\end{ex}

% nguồn: Ok_TracNghiemLop11_Lan1
\begin{ex}%[1D6H1-4]%[ma:974cb096]
	Cho $a=3^{\sqrt{5}}, b=3^2$ và $c=3^{\sqrt{6}}$. Mệnh đề nào dưới  đây đúng?
	\choice 
		 {$a<b<c$}
		{$a<c<b$} 
		{$c<a<b$} 
		 {\True $b<a<c$}
	 \loigiai{
		Do $2<\sqrt{5}<\sqrt{6}$ và $3>1$  nên
		 $$
		3^2 < 3^{\sqrt{5}}  < 3^{\sqrt{6}}
		\Rightarrow 
		b < a < c. 
		 $$
	} 
\end{ex}

% nguồn: Ok_TracNghiemLop11_Lan1
\begin{ex}%[1D6H1-4]%[ma:27c255cd]
	Nếu $a^{\frac{1}{3}}>a^{\frac{1}{6}}$ và $b^{\sqrt 3}>b^{\sqrt 5}$ thì 
	\choice 
		{$  a<1;0<b<1$}
		{$  a>1;b<1$}
		{$0<a<1;b<1$} 
		{\True $  a>1;0<b<1$}
	\loigiai{ 
		Ta có: $\dfrac{1}{3}>\dfrac{1}{6}$, lại  có $a^{\frac{1}{3}}>a^{\frac{1}{6}}\Rightarrow a>1$.\\
		Ta có:  $\sqrt 3<\sqrt 5$, lại có $b^{\sqrt 3}>b^{\sqrt 5}\Rightarrow 0<b<1$.
	} 
\end{ex}

% nguồn: Ok_TracNghiemLop11_Lan1
\begin{ex}%[1D6H1-4]%[ma:befd7691]
	Khẳng định nào dưới đây là  đúng?
	 \choice
		{$\left(\dfrac{1}{4}\right)^{-50}  < (\sqrt{2})^{100}$}
		{$3^{-\sqrt{2}} <  \left(\dfrac{1}{5}\right)^{\sqrt{2}}$}
		{$\left(\dfrac{3}{7}\right)^{\sqrt{3}} > \left(\dfrac{5}{8}\right)^{\sqrt{3}}$} 
		{\True $\left(\dfrac{1}{2}\right)^{-\pi} < \left(\dfrac{1}{3}\right)^{-\pi}$} 
	 \loigiai{
		 Ta có $\left(\dfrac{1}{2}\right)^{-\pi} = 2^{\pi}$ và $ \left(\dfrac{1}{3}\right)^{-\pi} = 3^{\pi}$.\\
		Vì $2^{\pi}  < 3^{\pi}$ nên $\left(\dfrac{1}{2}\right)^{-\pi} < \left(\dfrac{1}{3}\right)^{-\pi}$.
	 }
\end{ex}

% nguồn: Ok_TracNghiemLop11_Lan1
\begin{ex}%[1D6H1-4]%[ma:b8d35832]
	So sánh ba số: $\left(0,2\right)^{0,3}$, $\left(0,7\right)^{3,2}$ và  $\left(\sqrt 3\right) ^{0,2}$ ta được
	 \choice
		{\True $\left(0,7\right)^{3,2}<\left(0,2\right)^{0,3}<\left(\sqrt  3\right)^{0,2}$}
		{$\left(0,2\right)^{0,3}<\left(0,7\right)^{3,2}<\left(\sqrt 3\right)^{0,2}$} 
		{$\left(\sqrt 3\right)  ^{0,2}<\left(0,2\right)^{0,3}<\left(0,7\right)^{3,2}$}
		{$\left(0,2\right)^{0,3}<\left(\sqrt  3\right)^{0,2}<\left(0,7\right)^{3,2}$}
	\loigiai{ 
		Ta  có $\left(0,2\right)^{0,3}=\left(0,2\right)^{\tfrac{3}{10}}=\left[\left(0,2\right)^3\right]^{\tfrac{1}{10}}=\left(0,008\right)^{\tfrac{1}{10}}$.\\
		Mặt  khác
		$\left(0,7\right)^{3,2}=\left(0,7\right)^{\tfrac{32}{10}}=\left[\left(0,7\right)^{32}\right]^{\tfrac{1}{10}}$ 
		 và
		$\left(\sqrt  3\right) ^{0,2}=(3)^{\tfrac{1}{2}.\tfrac{2}{10}}=3^{\tfrac{1}{10}}$.\\
	Do  $\left(0,7\right)^{32}< 0,008 < 3$ nên $\left(0,7\right)^{3,2}<\left(0,2\right)^{0,3}<\left(\sqrt 3\right) ^{0,2}$.}
\end{ex}

% nguồn: Ok_TracNghiemLop11_Lan1
\begin{ex}%[1D6H1-4]%[ma:3b63344f]
	Biết rằng $2^{\alpha}=\dfrac{1}{4}$ và $3^{\alpha}=\dfrac{1}{9}$.  Khẳng định nào đúng?
	 \choice
		{\True  $6^{-2\alpha}<6^5$}
		{$6^{-2\alpha}>6^4$} 
		 {$6^{-2\alpha}<6^3$}
		 {$6^{-2\alpha}>6^5$}
	 \loigiai{
		Ta có $6^{\alpha}=\left(2\cdot  3\right)^{\alpha}=2^{\alpha}\cdot 3^{\alpha}=\dfrac{1}{4}\cdot \dfrac{1}{9}=\dfrac{1}{36}\Rightarrow 6^{-\alpha}=36\Rightarrow 6^{-2\alpha}=36^2=6^4$.\\
		Vậy $6^{-2\alpha}>6^4$  và $6^{-2\alpha}<6^3$ là sai.\\
		 Mà $6^{-2\alpha}=6^4<6^5$ nên A đúng.
	 }
\end{ex}

% nguồn: Ok_TracNghiemLop11_Lan1
\begin{ex}%[1D6H1-4]%[ma:17fc6b38]
	 Cho $ a>1$. Mệnh đề nào đúng?
	\choice 
		 {$\dfrac{1}{a^{2016}}<\dfrac{1}{a^{2017}}$}
		{$a^{\frac{1}{3}}>\sqrt  {a}$}
		{\True $a^{-\sqrt  3}>\dfrac{1}{a^{\sqrt 5}}$}
		{$\dfrac{\sqrt [3]{a^2}}{a}>1$} 
	\loigiai{ 
		 Vì $ a>1$ nên ${a^{\sqrt 3}}<a^{\sqrt 5}\Rightarrow \dfrac{1}{a^{\sqrt 3}}>\dfrac{1}{a^{\sqrt 5}}\Rightarrow a^{-\sqrt 3}>\dfrac{1}{a^{\sqrt 5}}$.
	} 
\end{ex}

% nguồn: Ok_TracNghiemLop11_Lan1
\begin{ex}%[1D6H1-4]%[ma:3e9cb46c]
	 Mệnh đề nào sai?
	\choice 
		{\True  $\left(\sqrt 3-1\right)^{2018}>\left(\sqrt 3-1\right)^{2017}$}
		{$2^{\sqrt 2+1}>2^{\sqrt 3}$} 
		{$\left(\sqrt  2-1\right)^{2017}>\left(\sqrt 2-1\right)^{2018}$}
		{$\left(1-\dfrac{\sqrt 2}{2}\right)^{2019}<\left(1-\dfrac{\sqrt  2}{2}\right)^{2018}$}
	\loigiai{ 
		\begin{itemize}
			\item $\left(\sqrt  3-1\right)^{2018}>\left(\sqrt 3-1\right)^{2017}$. Cùng cơ số, $0<\sqrt 3-1<1$, hàm nghịch biến, số mũ lớn hơn nên bé hơn nên A sai.
			\item $2^{\sqrt 2+1}>2^{\sqrt 3}$. Cùng cơ số, $2>1$, hàm đồng biến, số mũ $\left(\sqrt 2+1\right)^2=3+2\sqrt 2>\left(\sqrt 3\right)^2=3$ nên  lớn hơn nên B đúng.
			\item $\left(\sqrt 2-1\right)^{2017}>\left(\sqrt 2-1\right)^{2018}$. Cùng cơ số, $0<\sqrt 2-1<1$,  hàm nghịch biến, số mũ bé hơn nên lớn hơn nên C đúng.
			 \item $\left(1-\dfrac{\sqrt 2}{2}\right)^{2019}<\left(1-\dfrac{\sqrt 2}{2}\right)^{2018}$. Cùng cơ số, $0<1-\dfrac{\sqrt 2}{2}<1$, hàm nghịch biến, số mũ lớn hơn nên bé hơn nên D đúng.
		\end{itemize}
	 }
\end{ex}

% nguồn: Ok_TracNghiemLop11_Lan1
\begin{ex}%[1D6H1-4]%[ma:d8677b1f]
	Tìm tập tất cả các giá trị của $a$ để $\sqrt[{21}]{a^5}  > \sqrt[{7}]{a^2}$.
	\choice 
		{$a  > 0$}
		{\True$0 < a  < 1$}
		 {$a > 1$}
		{$\dfrac{5}{21}  < a < \dfrac{2}{7}$}
	\loigiai{ 
		Ta  có $\sqrt[{21}]{a^5} > \sqrt[{7}]{a^2} \Leftrightarrow \sqrt[{21}]{a^5} > \sqrt[{21}]{a^6}$ mà $5 < 6$. Vậy $0 < a < 1$.
	 }
\end{ex}

% nguồn: Ok_TracNghiemLop11_Lan1
\begin{ex}%[1D6H1-4]%[ma:2a34954e]
	 Khẳng định nào đúng?
	 \choice
		 {$\left(\dfrac{3}{7}\right)^{\sqrt 3}>\left(\dfrac{5}{8}\right)^{\sqrt 3}$}
		 {\True $\left(\dfrac{1}{2}\right)^{-\pi}<\left(\dfrac{1}{3}\right)^{-\pi}$}
		{$3^{-\sqrt 2}<\left(\dfrac{1}{5}\right)^{\sqrt 2}$} 
		 {$\left(\dfrac{1}{4}\right)^{-50}<\left(\sqrt 2\right)^{100}$}
	\loigiai{ 
		Ta có: 
		\begin{itemize}
			\item $\left(\dfrac{3}{7}\right)<\left(\dfrac{5}{8}\right)\Rightarrow \left(\dfrac{3}{7}\right)^{\sqrt 3}<\left(\dfrac{5}{8}\right)^{\sqrt 3}$. Phương án A  sai.
			 \item $\dfrac{1}{2}>\dfrac{1}{3}\Rightarrow \left(\dfrac{1}{2}\right)^{-\pi}<\left(\dfrac{1}{3}\right)^{-\pi}$. Phương án B đúng.
			\item $3<5\Rightarrow 3^{-\sqrt 2}>5^{-\sqrt 2}\Rightarrow 3^{-\sqrt  2}>\left(\dfrac{1}{5}\right)^{\sqrt 2}$. Phương án C sai.
			\item $\left(\dfrac{1}{4}\right)^{-50}<\left(\sqrt 2\right)^{100}\Rightarrow \left(2^{-2}\right)^{-50}<\left(2\right)^{100}\Rightarrow 2^{100}<2^{100}$. Phương án D sai. 
		\end{itemize}
	} 
\end{ex}

% nguồn: Ok_TracNghiemLop11_Lan1
\begin{ex}%[1D6H1-4]%[ma:6feeb633]
	Cho $\left(\sqrt{2}-1\right)^m < \left(\sqrt{2}-1\right)^n$. Khi  đó
	 \choice
		{$m=n$} 
		 {$m < n$}
		{\True  $m > n$}
		 {$m\ne n$}
	 \loigiai{
		Do $0<  \sqrt{2}-1< 1$ nên $\left(\sqrt{2}-1\right)^m < \left(\sqrt{2}-1\right)^n \Leftrightarrow m > n$.
	 }
\end{ex}


%%%%%%%%%%%%%%%% ĐÚNG/SAI %%%%%%%%%%%%%%%%

% nguồn: Ok_dungsai_lop11_Lan1
\begin{ex}%[1D6N1-4]%[ma:af0662d5]
 Với $a>0$, $b>0$, $m$, $n$ là các số thực tùy ý. Giả  sử các biểu thức xuất hiện trong các công thức của mỗi mệnh đề đều có nghĩa.
  \choiceTF
   {\True $a^m\cdot a^n=a^{m+n}$}
  {\True$a^{-n}=\dfrac{1}{a^n}$ với $n$  nguyên dương}
  {\True Nếu $\left(\dfrac{5}{2}\right)^m>\left(\dfrac{5}{2}\right)^n$  thì $m>n$}
  {Nếu $\left(\dfrac{\pi}{4}\right)^m>\left(\dfrac{\pi}{4}\right)^{3m-2}$ thì  $m<1$}
  \loigiai{
  \begin{itemchoice}
   \itemch  \textbf{Đúng}\\
   Ta có $a^m\cdot a^n=a^{m+n}$ với $a>0$ và $m$, $n$  là các số thực tùy ý.
   \itemch \textbf{Đúng}\\ 
   Ta có $a^{-n}=\dfrac{1}{a^n}$ với $n$  nguyên dương.
   \itemch \textbf{Đúng}\\ 
   Vì $\dfrac{5}{2}>1$  nên nếu $\left(\dfrac{5}{2}\right)^m>\left(\dfrac{5}{2}\right)^n$ thì $m>n$.
   \itemch \textbf{Sai} \\ 
   Vì  $0<\dfrac{\pi}{4}<1$ nên
   \[\left(\dfrac{\pi}{4}\right)^m>\left(\dfrac{\pi}{4}\right)^{3m-2}\Leftrightarrow m<3m-2\Leftrightarrow m>1.\] 
  \end{itemchoice}
 } 
\end{ex}


%%%%%%%%%%%%%%%% TỰ LUẬN %%%%%%%%%%%%%%%%

% nguồn: Ok_TuanLuan_Toan11_Lan1
\begin{bt}%[1D6V1-4]%[ma:bbee9bd2]
	Thực hiện các yêu  cầu sau:
	\begin{enumerate}
		\item Cho $A=199^{2023}$, $B=\sqrt{199}^{2024}$. So sánh  $A$, $B$.
		\item Sắp theo $A=3^{4999}$, $B=11^{4001}$ và $C=1331^{1000}$  theo thứ tự từ lớn đến bé.
	\end{enumerate}
	\loigiai{ 
		\begin{enumerate}
			\item Ta có  $199> 1$ nên $199^{2023} > 199^{1012} \Rightarrow A> B$. \\
			Vậy  $A> B$.
			\item 
			\begin{itemize}
				\item $C=1331^{1000}=\left(11^3 \right)^{1000}=11^{3000}$. 
				\item $11> 1$; $11^{4001} > 11^{2000} \Rightarrow  B> C$.
				\item $3^{4999} < 3^{5000}=\left(3^5 \right)^{1000}=243^{1000} < 1331^{1000}\Rightarrow  A< C$.
			\end{itemize}
			Vậy $B> C>  A$.
		\end{enumerate}
	 }
\end{bt}

% nguồn: Ok_TuanLuan_Toan11_Lan1
\begin{bt}%[1D6V1-4]%[ma:f4549564]
	Với những  giá trị nào của $a$ thì:
	\begin{enumerate}
		\item $a^\mathrm{e} >  a^{\pi}$;
		\item  $(a-1)^2 > (a-1)^{\sqrt{5}}$.
	\end{enumerate}
	 \loigiai{
		\begin{enumerate}
			 \item $a^\mathrm{e} > a^{\pi}$.\\
			Điều kiện: $a > 0$.\\ 
			Vì $\heva{&\mathrm{e} < \pi\\&a^\mathrm{e} > a^{\pi}}\Rightarrow 0< a <  1$.\\
			Vậy $0< a <  1$.
			\item $(a-1)^2 > (a-1)^{\sqrt{5}}$.\\ 
			Điều kiện:  $a-1> 0\Leftrightarrow a > 1$.\\
			Vì  $\heva{&2< \sqrt{5}\\
			&\left(a-1\right)^2>  \left(a-1\right)^{\sqrt{5}}}\Rightarrow 0< a-1< 1\Leftrightarrow 1< a < 2$.\\
			Vậy $1< a < 2$. 
		\end{enumerate}
	 }
\end{bt}

% nguồn: Ok_TuanLuan_Toan11_Lan2
\begin{ex}%[1D6V1-4]%[ma:d38389d0]
	Thực hiện các  yêu cầu sau:
	\begin{enumerate}
		\item  Cho $A=199^{2023}$, $B=\sqrt{199}^{2024}$. So sánh $A$, $B$.
		\item Sắp theo $A=3^{4999}$, $B=11^{4001}$ và $C=1331^{1000}$ theo thứ tự từ lớn đến bé. 
	\end{enumerate}
	\loigiai{ 
		\begin{enumerate}
			\item  Ta có $199> 1$ nên $199^{2023} > 199^{1012} \Rightarrow A> B$. \\
			 Vậy $A> B$.
			 \item
			\begin{itemize}
				\item $C=1331^{1000}=\left(11^3  \right)^{1000}=11^{3000}$.
				\item  $11> 1$; $11^{4001} > 11^{2000} \Rightarrow B> C$.
				\item $3^{4999} < 3^{5000}=\left(3^5 \right)^{1000}=243^{1000} <  1331^{1000}\Rightarrow A< C$.
			\end{itemize}
			Vậy $B>  C> A$.
		\end{enumerate}
	} 
\end{ex}

% nguồn: Ok_TuanLuan_Toan11_Lan2
\begin{ex}%[1D6V1-4]%[ma:bcb8d287]
	Với những giá trị nào của $a$ thì: 
	\begin{enumerate}
		 \item $a^\mathrm{e} > a^{\pi}$;
		 \item $(a-1)^2 > (a-1)^{\sqrt{5}}$.
	\end{enumerate}
	\loigiai{ 
		\begin{enumerate}
			\item $a^\mathrm{e} >  a^{\pi}$.\\
			Điều kiện: $a >  0$.\\
			Vì  $\heva{&\mathrm{e} < \pi\\&a^\mathrm{e} > a^{\pi}}\Rightarrow 0< a < 1$.\\
			Vậy $0< a  < 1$.
			\item $(a-1)^2  > (a-1)^{\sqrt{5}}$.\\
			Điều kiện: $a-1> 0\Leftrightarrow  a > 1$.\\
			 Vì $\heva{&2< \sqrt{5}\\
			&\left(a-1\right)^2> \left(a-1\right)^{\sqrt{5}}}\Rightarrow 0<  a-1< 1\Leftrightarrow 1< a < 2$.\\
			 Vậy $1< a < 2$.
		\end{enumerate}
	 }
\end{ex}
